Ten young adults crossed obstacles of three heights (10, 20 and 30\,\% of leg length) and two widths (5 and 15\,cm) while wearing
six Movella Dot IMUs (thigh, shank and foot of both legs), with camera motion capture as reference. From the processed recordings
a crossing-level dataset was built (349 rows: 303 obstacle crossings, 46 level walks; 260 crossings used for machine learning), and
nine MATLAB analyses were run, every accuracy being measured on a participant left out of training. People lift each foot by about
the obstacle height (10.1\,mm per cm for the leading foot, 95\,\% CI 9.4--10.8) and keep a constant safety margin, but each person
has a strong walking signature and a personal obstacle margin. Generic models therefore recognise the obstacle height of a new
person in only 56--73\,\% of crossings (chance 33\,\%), whereas scaling the features with the person's own crossings, without any
label, gives 90--92\,\%; two foot IMUs are enough (92.3\,\%, 95\,\% CI 87--96), the height can be estimated in centimetres with a
2.4\,cm error, and a camera-free pipeline that finds the crossings, pairs the two limbs and classifies the height is fully correct
for 86\,\% of the crossings. Obstacle width is not recognisable (54--66\,\%, chance 50\,\%) because people barely change how they
cross a deeper obstacle. For a haptic cue to the trailing limb, whose deadline is the moment its toe leaves the floor, the information must come from the
leading limb: the trailing limb shows no sign of the obstacle height before it lifts, the leading-foot apex leaves a median 0.45\,s
and the leading landing only 0.2\,s, and a single leading-limb signal (foot height, or thigh flexion when only angle sensors are used)
detects the crossing in 97--99\,\% of cases and grades the height correctly in 64--73\,\% of crossings with personal thresholds.

